\chapter{Configuration}
\label{ch:configuration}

\section{Configuration surfaces}

Expo Snips has no configuration file of its own and no server-side
configuration. Configuration exists at four points:

\begin{enumerate}
  \item \textbf{CLI flags and environment variables} (command line / process
        environment).
  \item \textbf{VS Code settings} (editor settings, per user or workspace).
  \item \textbf{Recipe definitions} (JSON data that configures generation
        behaviour per screen).
  \item \textbf{Per-project state} (the \texttt{.expo-snips.json} manifest and
        the detected stack of the user's project).
\end{enumerate}

There is no \texttt{.env} file in the repository and none is read by the code.

\section{Environment variables}

\begin{longtable}{p{4.4cm}p{2.2cm}p{7.6cm}}
\caption{Environment variables}\\
\toprule
\textbf{Variable} & \textbf{Default} & \textbf{Effect} \\
\midrule

\texttt{EXPO\_SNIPS\_REGISTRY} & empty (embedded catalog) & Overrides the
registry base URL used by the CLI. An empty value selects the embedded catalog. \\
\texttt{PORT} & \texttt{3000} & Port for the development registry server
(\texttt{scripts/serve.mjs}). \\
\texttt{NPM\_TOKEN} & --- & CI only. Authenticates npm publishing in
\texttt{publish.yml}. Never read by application code. \\
\texttt{GITHUB\_TOKEN} & --- & CI only. Used by the Pages deploy action. Never
read by application code. \\
\bottomrule
\end{longtable}

\textbf{Precedence (CLI):} an explicit \texttt{--registry} flag takes precedence
over \texttt{EXPO\_SNIPS\_REGISTRY}; when neither is set, the built-in default
registry constant applies. An empty string in either place means ``use the
catalog bundled with the package''.

\section{CLI options}

Global and per-command options are defined with \texttt{commander}. The
generation-relevant options are:

\begin{longtable}{p{3.4cm}p{3.4cm}p{6.8cm}}
\caption{CLI generation options}\\
\toprule
\textbf{Option} & \textbf{Values} & \textbf{Effect} \\
\midrule

\texttt{--styling} & \texttt{nativewind}, \texttt{stylesheet},
\texttt{tamagui} & Selects the styling template branch; validated against the
recipe's \texttt{supported.styling}. \\
\texttt{--router} & \texttt{expo-router}, \texttt{react-navigation} & Selects the
navigation template branch; validated against
\texttt{supported.router}. \\
\texttt{--ui} & \texttt{native}, \texttt{shadcn}, \texttt{custom} & Selects the UI
primitive branch; \texttt{shadcn} additionally forces
\texttt{styling=nativewind} and \texttt{icons=lucide} and pulls in primitive
recipes. \\
\texttt{--icons} & \texttt{lucide}, \texttt{vector-icons}, \texttt{none} &
Selects the icon template branch. \\
\texttt{--variant} & \textit{recipe-defined} & Selects a named template variant.
Not validated against the recipe. \\
\texttt{--typescript} / \texttt{--javascript} & flag & Forces the language branch
(overrides detection). \\
\texttt{--no-install} & flag & Skips dependency installation. \\
\texttt{--registry <url>} & URL & Overrides the registry base URL. \\
\bottomrule
\end{longtable}

When an option is omitted, the value is inferred from the project (see below),
then falls back to a default, then to an interactive prompt if the terminal is
interactive.

\section{VS Code settings}

The extension contributes a configuration section titled ``Expo Snips'' with
these properties:

\begin{longtable}{p{4.2cm}p{1.8cm}p{2.4cm}p{5.8cm}}
\caption{VS Code settings}\\
\toprule
\textbf{Setting} & \textbf{Type} & \textbf{Default} & \textbf{Effect} \\
\midrule

\texttt{expo-snips.registryUrl} & string & \textit{empty} & Remote registry base.
Empty uses the embedded catalog. \\
\texttt{expo-snips.defaultStyling} & enum & \texttt{nativewind} & Fallback styling
when detection fails. Values: \texttt{nativewind}, \texttt{stylesheet},
\texttt{tamagui}. \\
\texttt{expo-snips.defaultRouter} & enum & \texttt{expo-router} & Fallback router.
Values: \texttt{expo-router}, \texttt{react-navigation}. \\
\texttt{expo-snips.defaultUi} & enum & \texttt{native} & Fallback UI primitives.
Values: \texttt{native}, \texttt{shadcn}, \texttt{custom}. \\
\texttt{expo-snips.defaultIcons} & enum & \texttt{lucide} & Fallback icon library.
Values: \texttt{lucide}, \texttt{vector-icons}, \texttt{none}. \\
\texttt{expo-snips.installDependencies} & boolean & \texttt{true} & Install
recipe dependencies in a terminal after generation. \\
\bottomrule
\end{longtable}

\textbf{Configuration mismatch to note:} the settings enums advertise
\texttt{tamagui}, \texttt{vector-icons} and \texttt{custom}, but no recipe in the
catalog currently implements them (every recipe's \texttt{supported} array
excludes them). Selecting such a fallback would have no recipe to match. The
extension README acknowledges this. This is a documentation/configuration
inconsistency rather than a runtime defect, but it can mislead users.

\section{Project detection (configuration by inference)}

The tool configures itself from the target project rather than requiring input.
Detection reads:

\begin{longtable}{p{4.6cm}p{9.6cm}}
\caption{Detection probes}\\
\toprule
\textbf{Setting} & \textbf{Probe} \\
\midrule

Router & \texttt{expo-router} present in dependencies selects
\texttt{expo-router}; otherwise \texttt{@react-navigation/native} selects
\texttt{react-navigation}. \\
Styling & \texttt{nativewind} in dependencies selects \texttt{nativewind};
otherwise \texttt{stylesheet}. \\
UI primitives & \texttt{@shadcn/ui} or \texttt{react-native-reusables} selects
\texttt{shadcn}; otherwise \texttt{native}. \\
Icons & \texttt{lucide-react-native} selects \texttt{lucide}; otherwise
\texttt{none}. \\
TypeScript & Presence of \texttt{tsconfig.json}. \\
\bottomrule
\end{longtable}

\textbf{Status:} IMPLEMENTED. Detected values are used as defaults; explicit
flags (CLI) or settings fallbacks (extension) take precedence.

\section{Recipe-level configuration}

Recipes are the richest configuration surface. A recipe declares:

\begin{itemize}
  \item its identity and version (\texttt{id}, \texttt{name}, \texttt{description},
        \texttt{version});
  \item the options it \emph{supports} (\texttt{supported.styling/router/ui/icons})
        --- the tool will only generate combinations a recipe declares;
  \item its \texttt{dependencies} and \texttt{devDependencies};
  \item its \texttt{files}, each with a target \texttt{path}, a
        \texttt{templateFile}, an optional \texttt{when} condition, and optional
        \texttt{variants} and \texttt{perStyling} template overrides;
  \item optional \texttt{variants} (named style alternatives),
        \texttt{tags} (search keywords) and \texttt{primitives} (UI primitive
        recipes required when \texttt{ui=shadcn}).
\end{itemize}

The authoritative machine-readable definition is
\texttt{packages/registry/schema.json} (JSON Schema Draft-07), reproduced in
Chapter~\ref{ch:registry-and-recipes}.

\section{Project state: the manifest}

\begin{lstlisting}[language=JavaScript]
// .expo-snips.json  (written into the user's project root)
interface Manifest {
  version: number;
  screens: Array<{
    recipeId: string;
    recipeVersion: string;
    variant?: string;
    options: GenerationOptions;
    files: string[];
    installedAt: string;
  }>;
}
\end{lstlisting}

The manifest is created by the tool, written with two-space indentation, and
tolerates absent or unreadable content by treating the project as having no
tracked screens. The upsert key for a screen is the recipe id \emph{plus} the
variant, so the same recipe can be tracked once per variant.

\section{Build and development configuration}

\begin{longtable}{p{4.4cm}p{9.8cm}}
\caption{Build-related configuration}\\
\toprule
\textbf{File} & \textbf{Purpose} \\
\midrule

\texttt{pnpm-workspace.yaml} & Declares \texttt{packages/*} and an
\texttt{allowBuilds} policy that enables build scripts for \texttt{esbuild} and
\texttt{@vscode/vsce-sign} and disables them for \texttt{keytar}. \\
\texttt{tsconfig.base.json} & Shared TypeScript compiler options extended by each
package. \\
Per-package \texttt{tsconfig.json} & Package-specific compilation settings. \\
\texttt{esbuild.mjs} & Bundle configuration for the extension (CJS, Node 18
target, \texttt{vscode} external, sourcemap on). \\
\texttt{.vscodeignore} & Excludes sources and dev material from the packaged
VSIX. \\
\texttt{.gitignore} & Excludes build output and generated files from version
control. \\
\bottomrule
\end{longtable}

There is no \texttt{eslint}, \texttt{prettier} or \texttt{editorconfig}
configuration in the repository.

\section{Configuration summary}

\begin{table}[H]
\centering
\caption{Configuration posture}
\begin{tabular}{p{4.4cm}p{9.8cm}}
\toprule
\textbf{Surface} & \textbf{Assessment} \\
\midrule
CLI flags / env & IMPLEMENTED; documented in Chapter~\ref{ch:cli-reference}. \\
Editor settings & IMPLEMENTED; three values advertise unimplemented options. \\
Recipe data & IMPLEMENTED; schema-governed. \\
Detection & IMPLEMENTED; four probes plus a TypeScript probe. \\
Config files & NOT APPLICABLE (none exists). \\
Secrets & None in the application; CI secrets only. \\
\bottomrule
\end{tabular}
\end{table}
